\section{Interfejs silnika}
\label{sec:interfejs-silnika}

Silnik domenowy udostępnia niewielki interfejs publiczny, za pomocą którego
można rozpocząć rozdanie, pobrać obserwację, wykonać decyzję oraz odczytać
wynik zakończonego epizodu. Interfejs nie zależy od konkretnego algorytmu
uczenia ani od biblioteki Gymnasium.

Centralnym elementem jest klasa \texttt{UTHGame}. Pojedyncza instancja
przechowuje stan aktualnego rozdania, źródło kart oraz opcjonalną historię
decyzji. Po zakończeniu rozdania ta sama instancja może zostać wykorzystana
ponownie przez wywołanie metody \texttt{reset()}.

\subsection{Tworzenie silnika}
\label{subsec:tworzenie-silnika}

Konstruktor klasy \texttt{UTHGame} przyjmuje dwa opcjonalne parametry:

\begin{itemize}
    \item \texttt{seed} -- ziarno nadrzędnego generatora liczb
          pseudolosowych,
    \item \texttt{record\_history} -- informację, czy należy rejestrować
          przebieg kolejnych decyzji.
\end{itemize}

Ziarno określa powtarzalny strumień talii generowanych dla kolejnych
rozdawań. Parametr nie ustala pojedynczej, niezmiennej talii dla całej
instancji. Każde wywołanie \texttt{reset()} tworzy nową talię, której ziarno
jest pobierane z generatora środowiska.

Rejestrowanie historii jest domyślnie wyłączone. Dzięki temu podstawowe
wykonywanie dużej liczby rozdań nie wymaga przechowywania dodatkowych
obiektów diagnostycznych.

\subsection{Rozpoczęcie epizodu}
\label{subsec:interfejs-reset}

Metoda \texttt{reset()} rozpoczyna nowe rozdanie i zwraca pierwszą obserwację
w fazie \texttt{PREFLOP}. W standardowym użyciu środowisko samodzielnie
tworzy przetasowaną talię. Metoda umożliwia również przekazanie własnego
obiektu implementującego interfejs \texttt{CardSource}, na przykład
deterministycznej talii \texttt{FixedDeck}.

Sygnaturę operacji można przedstawić jako

\begin{equation}
    reset :
    CardSource \cup \{None\}
    \rightarrow
    UTHObservation.
\end{equation}

Rozpoczęcie rozdania obejmuje:

\begin{enumerate}
    \item wybór lub utworzenie źródła kart,
    \item rozdanie dwóch kart graczowi i dwóch kart krupierowi,
    \item utworzenie stanu \texttt{PREFLOP},
    \item wyzerowanie historii decyzji bieżącego rozdania,
    \item utworzenie obserwacji gracza.
\end{enumerate}

Reset jest wykonywany transakcyjnie. Silnik najpierw próbuje pobrać wszystkie
karty potrzebne do utworzenia stanu początkowego, a dopiero po powodzeniu
zapisuje nowy stan i źródło kart. Jeżeli przekazane źródło nie może dostarczyć
wymaganych kart, poprzedni stan silnika nie zostaje zastąpiony stanem
częściowo utworzonego rozdania.

\subsection{Wykonanie decyzji}
\label{subsec:interfejs-step}

Metoda \texttt{step()} przyjmuje jedną wartość typu \texttt{Action}
i wykonuje odpowiadające jej przejście. Jej uproszczona sygnatura ma postać

\begin{equation}
    step :
    Action
    \rightarrow
    StepResult.
\end{equation}

Wynikiem operacji jest obiekt \texttt{StepResult}, zawierający:

\begin{itemize}
    \item obserwację po wykonaniu akcji,
    \item wynik rozdania, jeżeli osiągnięto stan terminalny,
    \item rozliczenie zakładów w stanie terminalnym,
    \item szczegóły showdownu, jeżeli rozdanie nie zakończyło się foldem.
\end{itemize}

Wynik kroku pośredniego zawiera jedynie kolejną obserwację. Pola dotyczące
wyniku i rozliczenia pozostają wtedy nieokreślone. W kroku terminalnym
obiekt zawiera komplet informacji niezbędnych do obliczenia nagrody
i zarejestrowania wyniku eksperymentu.

Właściwość \texttt{terminated} obiektu \texttt{StepResult} informuje,
czy wykonana akcja zakończyła rozdanie. Jej wartość jest wyznaczana na
podstawie fazy zawartej w obserwacji.

\subsection{Właściwości publiczne}
\label{subsec:wlasciwosci-publiczne}

Poza metodami \texttt{reset()} i \texttt{step()} silnik udostępnia
właściwości przedstawione w tabeli~\ref{tab:interfejs-uth-game}.

\begin{table}[htbp]
    \centering
    \caption{Publiczny interfejs klasy \texttt{UTHGame}}
    \label{tab:interfejs-uth-game}
    \begin{tabular}{
        |p{0.25\textwidth}
        |p{0.20\textwidth}
        |p{0.45\textwidth}|
    }
        \hline
        \textbf{Element}
        &
        \textbf{Typ wyniku}
        &
        \textbf{Znaczenie}
        \\
        \hline

        \texttt{reset()}
        &
        \texttt{UTHObservation}
        &
        Rozpoczyna nowe rozdanie i zwraca obserwację preflop.
        \\
        \hline

        \texttt{step(action)}
        &
        \texttt{StepResult}
        &
        Wykonuje legalną akcję i zwraca wynik przejścia.
        \\
        \hline

        \texttt{observation}
        &
        \texttt{UTHObservation}
        &
        Zwraca aktualne informacje widoczne dla agenta.
        \\
        \hline

        \texttt{state}
        &
        \texttt{RoundState}
        &
        Zwraca pełny stan wewnętrzny, przeznaczony głównie do testów
        i diagnostyki.
        \\
        \hline

        \texttt{is\_started}
        &
        \texttt{bool}
        &
        Informuje, czy rozpoczęto co najmniej jedno rozdanie.
        \\
        \hline

        \texttt{history\_enabled}
        &
        \texttt{bool}
        &
        Informuje, czy rejestrowanie historii jest aktywne.
        \\
        \hline

        \texttt{trace}
        &
        \texttt{RoundTrace} lub \texttt{None}
        &
        Zwraca zapis bieżącego rozdania, jeżeli historia jest włączona.
        \\
        \hline
    \end{tabular}
\end{table}

Właściwość \texttt{state} nie jest częścią interfejsu przeznaczonego
do podejmowania decyzji przez agenta. Udostępniono ją na potrzeby testów,
diagnostyki i analizy środowiska. Agent powinien otrzymywać wyłącznie
wartość właściwości \texttt{observation} albo obserwację zwróconą przez
\texttt{reset()} lub \texttt{step()}.
\subsection{Integracja z Gymnasium}
\label{subsec:integracja-gymnasium}

Silnik \texttt{UTHGame} jest niezależny od biblioteki Gymnasium. Integracja
z algorytmami uczenia przez wzmacnianie zostanie zrealizowana przez osobny
adapter, który nie będzie implementował zasad gry, lecz jedynie warstwę
translacyjną: numeryczne kodowanie obserwacji, odwzorowanie indeksów akcji
na wartości \texttt{Action}, utworzenie maski legalnych akcji oraz
wyznaczenie nagrody i flag \texttt{terminated} i \texttt{truncated}.
Wszystkie przejścia, walidacja akcji, showdown i rozliczenia pozostaną
odpowiedzialnością silnika domenowego, co zapobiega powstaniu dwóch
niezależnych implementacji UTH.

% TODO: uzupełnić szczegóły adaptera Gymnasium (przestrzenie akcji
% i obserwacji, słownik info, mapowanie reset/step) po jego implementacji.

\section{Rejestrowanie przebiegu i walidacja środowiska}
\label{sec:rejestrowanie-i-walidacja}

Poprawność środowiska ma bezpośredni wpływ na wyniki eksperymentów.
Błąd w kolejności rozdawania kart, rozliczaniu zakładów albo udostępnianiu
obserwacji mógłby zostać błędnie zinterpretowany jako zachowanie wyuczone
przez agenta. Z tego względu implementację uzupełniono o opcjonalny mechanizm
rejestrowania przebiegu rozdania oraz zestaw testów weryfikujących zarówno
pojedyncze komponenty, jak i pełne ścieżki rozgrywki.

\subsection{Opcjonalne rejestrowanie historii}
\label{subsec:opcjonalna-historia}

Rejestrowanie historii jest aktywowane parametrem
\texttt{record\_history} przekazywanym podczas tworzenia obiektu
\texttt{UTHGame}. Domyślnie mechanizm pozostaje wyłączony, aby standardowe
wykonywanie dużej liczby epizodów nie wymagało tworzenia dodatkowych
obiektów diagnostycznych.

Po włączeniu historii każde poprawnie rozpoczęte rozdanie otrzymuje dodatni,
rosnący identyfikator \texttt{round\_id}. Lista decyzji jest zerowana przy
każdym poprawnym wywołaniu metody \texttt{reset()}.

Właściwość \texttt{trace} zwraca obiekt \texttt{RoundTrace}. Jeżeli historia
jest wyłączona albo nie rozpoczęto jeszcze rozdania, jej wartością jest
\texttt{None}.

Mechanizm historii nie stanowi pamięci agenta i nie jest częścią obserwacji.
Jest przeznaczony do:

\begin{itemize}
    \item diagnozowania nieoczekiwanych decyzji,
    \item odtwarzania przebiegu wybranych rozdań,
    \item analizy działania agentów bazowych i uczących się,
    \item zapisywania przykładów błędów podczas eksperymentów,
    \item weryfikowania zgodności decyzji z legalną przestrzenią akcji.
\end{itemize}

\subsection{Rekord pojedynczej decyzji}
\label{subsec:decision-record}

Pojedynczą decyzję reprezentuje niemodyfikowalny obiekt
\texttt{DecisionRecord}. Zawiera on:

\begin{itemize}
    \item obserwację dostępną agentowi przed wykonaniem decyzji,
    \item wybraną akcję.
\end{itemize}

Rekord można przedstawić jako parę

\begin{equation}
    d_t =
    \left(
        O_t,
        A_t
    \right),
\end{equation}

gdzie \(O_t\) oznacza obserwację przed wykonaniem akcji, a \(A_t\) oznacza
decyzję agenta.

Zapisanie obserwacji sprzed przejścia jest istotne, ponieważ pozwala
odtworzyć informacje, na podstawie których agent rzeczywiście podjął
decyzję. Zapis obserwacji po wykonaniu akcji prowadziłby do niejednoznaczności:
mógłby zawierać karty odkryte dopiero jako skutek tej decyzji.

Podczas tworzenia rekordu sprawdzane jest, czy:

\begin{itemize}
    \item obserwacja jest obiektem \texttt{UTHObservation},
    \item akcja jest wartością typu \texttt{Action},
    \item obserwacja nie jest terminalna,
    \item akcja należy do zbioru legalnego dla zapisanej obserwacji.
\end{itemize}

Dzięki temu historia nie może zawierać decyzji wykonanej po zakończeniu
rozdania ani akcji sprzecznej z fazą zapisaną w obserwacji.

\subsection{Ślad całego rozdania}
\label{subsec:round-trace}

Obiekt \texttt{RoundTrace} reprezentuje niemodyfikowalny zapis aktualnego
przebiegu rozdania. Zawiera:

\begin{itemize}
    \item dodatni identyfikator rozdania,
    \item uporządkowaną krotkę obiektów \texttt{DecisionRecord},
    \item aktualny pełny stan \texttt{RoundState}.
\end{itemize}

Dla rozdania zawierającego \(n\) decyzji jego ślad można zapisać jako

\begin{equation}
    \tau =
    \left(
        id,
        d_0,
        d_1,
        \ldots,
        d_{n-1},
        S_n
    \right).
\end{equation}

Obserwacje znajdujące się w rekordach decyzji zawierają wyłącznie informacje
dostępne agentowi. Końcowy element \(S_n\) jest natomiast pełnym stanem
diagnostycznym i może zawierać karty krupiera, karty spalone oraz rozliczenie
rozdania.

Z tego powodu \texttt{RoundTrace} nie może być przekazywany agentowi jako
obserwacja. Jest to obiekt przeznaczony dla kodu eksperymentalnego, testów
i narzędzi analitycznych.

Ślad udostępnia również właściwości pomocnicze:

\begin{itemize}
    \item \texttt{completed} -- informację, czy osiągnięto stan terminalny,
    \item \texttt{outcome} -- końcowy wynik rozdania,
    \item \texttt{settlement} -- rozliczenie Ante, Blind i Play,
    \item \texttt{showdown} -- szczegóły porównania rąk, jeżeli do niego
          doszło.
\end{itemize}

Przed zakończeniem rozdania właściwości związane z wynikiem mogą zwracać
wartość \texttt{None}. Po foldzie dostępne jest rozliczenie i wynik, ale
nie występują szczegóły showdownu.

\subsection{Transakcyjne zapisywanie decyzji}
\label{subsec:transakcyjne-zapisywanie}

Metoda \texttt{step()} najpierw zapisuje bieżącą obserwację w zmiennej
lokalnej, następnie próbuje wykonać właściwe przejście, a dopiero po jego
powodzeniu dodaje rekord do historii.

Kolejność operacji ma postać

\begin{equation}
    O_t
    \rightarrow
    \operatorname{validate}(A_t)
    \rightarrow
    \operatorname{transition}(S_t,A_t)
    \rightarrow
    S_{t+1}
    \rightarrow
    \operatorname{record}(O_t,A_t).
\end{equation}

Jeżeli walidacja lub wykonywanie przejścia zakończy się wyjątkiem, etap
\(\operatorname{record}\) nie zostaje osiągnięty. Historia zawiera zatem
wyłącznie decyzje, dla których środowisko utworzyło poprawny kolejny stan.

Takie zachowanie zabezpiecza między innymi przed zapisaniem:

\begin{itemize}
    \item akcji nielegalnej w aktualnej fazie,
    \item decyzji po zakończeniu rozdania,
    \item akcji, podczas której źródło kart zwróciło niepoprawną wartość,
    \item niedokończonego przejścia wynikającego z pustej talii.
\end{itemize}

Ta sama zasada jest stosowana podczas resetowania środowiska. Nowy stan,
źródło kart, identyfikator rozdania i pusta historia są zapisywane dopiero
po poprawnym rozdaniu wszystkich czterech kart początkowych.

\subsection{Testy deterministyczne}
\label{subsec:testy-deterministyczne}

Znaczna część testów pełnych rozdań wykorzystuje deterministyczne źródło
\texttt{FixedDeck}. Jawne określenie kolejności kart pozwala przygotować
scenariusz, którego wynik nie zależy od generatora pseudolosowego.

Deterministyczny scenariusz może określać kolejno:

\begin{equation}
    P_1,
    D_1,
    P_2,
    D_2,
    B_1,
    F_1,
    F_2,
    F_3,
    B_2,
    T,
    R,
\end{equation}

gdzie \(P\) i \(D\) oznaczają karty własne gracza i krupiera, \(B\) karty
spalone, \(F\) karty flopa, \(T\) turn, a \(R\) river.

Tak przygotowane źródło umożliwia kontrolowanie jednocześnie:

\begin{itemize}
    \item kart obu stron,
    \item końcowego układu gracza i krupiera,
    \item kwalifikacji krupiera,
    \item wyniku showdownu,
    \item wypłaty zakładu Blind,
    \item łącznego wyniku netto.
\end{itemize}

W przeciwieństwie do testu opartego wyłącznie na ziarnie
\texttt{FixedDeck} jednoznacznie pokazuje w kodzie testowym, jakie karty
mają zostać dobrane. Ułatwia to analizę nieudanego przypadku i ogranicza
zależność testu od szczegółów algorytmu tasowania.

\subsection{Poziomy testowania}
\label{subsec:poziomy-testowania}

Testy podzielono według odpowiedzialności komponentów. Testy jednostkowe
sprawdzają w izolacji model kart i talię, evaluator, modele stanu,
rozdawanie, legalność akcji, showdown oraz rozliczenia. Testy integracyjne
wykonują kompletne rozdania obejmujące wszystkie ścieżki maszyny stanów, od
\texttt{reset()} do stanu terminalnego, wraz z rejestrowaniem historii. Samo
poprawne przetestowanie evaluatora i modułu rozliczeń nie gwarantuje bowiem,
że silnik przekaże im właściwe karty i mnożnik zakładu Play.

Poza przykładowymi wynikami testy weryfikują niezmienniki obowiązujące dla
wszystkich poprawnych rozdań, między innymi brak powtarzających się kart,
zgodność liczby kart wspólnych i spalonych z fazą, brak danych terminalnych
w stanie pośrednim, brak showdownu po spasowaniu, wykonanie co najwyżej
jednego zakładu Play, brak kart krupiera i kart spalonych w obserwacji oraz
brak zmiany stanu po odrzuconej akcji. Walidacja obejmuje zatem także
przypadki brzegowe i próby utworzenia stanów naruszających reguły domenowe.

\subsection{Pokrycie kodu}
\label{subsec:pokrycie-kodu}

Zestaw testów uruchamiany jest za pomocą biblioteki \texttt{pytest}.
Podczas prac nad silnikiem analizowano zarówno pokrycie instrukcji, jak
i pokrycie rozgałęzień. Pozwoliło to wykryć nieprzetestowane ścieżki
walidacji, przypadki brzegowe evaluatora oraz alternatywne sposoby
rozliczania zakładów.

Na etapie zakończenia implementacji podstawowego silnika uzyskano pełne
pokrycie instrukcji i rozgałęzień dla objętych pomiarem modułów. Wynik
pokrycia nie stanowi samodzielnego dowodu poprawności implementacji.
Informuje jedynie, że testy wykonały wszystkie mierzone linie i gałęzie
programu. Z tego względu został uzupełniony testami opartymi na konkretnych
scenariuszach domenowych oraz walidacją niezmienników.

\section{Ograniczenia środowiska}
\label{sec:ograniczenia-srodowiska}

Zaimplementowane środowisko odwzorowuje podstawowy wariant Ultimate Texas
Hold’em potrzebny do prowadzenia eksperymentów nad strategiami podejmowania
decyzji. Celem silnika nie jest pełna symulacja działalności stołu
kasynowego, lecz dostarczenie jednoznacznego, deterministycznie testowalnego
modelu procesu decyzyjnego gracza.

Zakres implementacji został celowo ograniczony do elementów wpływających
na decyzje Play i końcowy wynik podstawowych zakładów. Pozostałe mechanizmy
mogą zostać dodane w przyszłości jako osobne rozszerzenia, bez zmieniania
głównej maszyny stanów.

\subsection{Zakres podstawowego środowiska}
\label{subsec:zakres-podstawowego-srodowiska}

Zaimplementowany wariant środowiska obejmuje:

\begin{itemize}
    \item jednego gracza rozgrywającego rozdanie przeciwko krupierowi,
    \item standardową talię 52 kart bez jokerów,
    \item zakłady Ante, Blind i Play,
    \item decyzje Play wykonywane przed flopem, po flopie i po odkryciu
          całego stołu,
    \item kwalifikację krupiera od pary,
    \item standardową tabelę wypłat zakładu Blind,
    \item zakończenie przez showdown albo fold,
    \item dokładne rozliczenie finansowe w jednostkach Ante,
    \item kontrolowane źródło losowości i deterministyczne talie testowe,
    \item opcjonalne rejestrowanie przebiegu rozdania.
\end{itemize}

Poza jego zakresem pozostają natomiast zakłady boczne Trips i Progressive,
model wielu graczy przy stole, saldo gracza i limity stołu, interfejs
graficzny, a także adapter Gymnasium oraz numeryczne kodowanie obserwacji,
które zostaną dodane jako osobne warstwy.

\subsection{Pominięcie zakładów dodatkowych}
\label{subsec:pominiecie-zakladow-dodatkowych}

Zakład Trips jest rozliczany na podstawie końcowego układu gracza niezależnie
od wyniku porównania z krupierem. Nie wprowadza kolejnego momentu decyzyjnego
podczas rozdania, ponieważ jego wartość jest ustalana przed rozdaniem kart.
Wpływa jednak na rozkład końcowych wypłat.

Uwzględnienie Trips w podstawowej wersji środowiska utrudniałoby interpretację
wyników strategii Play. Agent mógłby uzyskiwać wysokie lub niskie wyniki
wynikające z zakładu, na który nie ma wpływu podczas właściwej sekwencji
decyzyjnej. Z tego względu bazowe eksperymenty będą oceniały wyłącznie wynik
Ante, Blind i Play.

Progressive również nie zmienia późniejszych decyzji Play, ale wymaga
dodatkowego modelu wypłat. W wariancie z narastającym jackpotem wartość
potencjalnej nagrody zależy ponadto od stanu utrzymywanego pomiędzy
rozdaniami. Oznaczałoby to, że epizody nie byłyby w pełni niezależne bez
dodatkowego modelowania wartości puli.

Oba zakłady mogą zostać w przyszłości dodane jako rozszerzenia warstwy
rozliczeń. Nie wymagają zmiany podstawowej kolejności faz
\texttt{PREFLOP}, \texttt{FLOP}, \texttt{RIVER} i \texttt{TERMINAL}, ale
zmieniają końcową funkcję wypłaty.

\subsection{Brak modelu salda i sesji}
\label{subsec:brak-modelu-salda}

Silnik traktuje każde rozdanie jako niezależny epizod. Nie przechowuje salda
gracza, historii wyników pomiędzy rozdaniami ani ograniczeń wynikających
z dostępnego kapitału.

Ante i Blind mają stałą wartość jednej jednostki, natomiast Play jest
wielokrotnością Ante wynikającą bezpośrednio z wybranej akcji. Dzięki temu
wyniki różnych strategii mogą być porównywane niezależnie od przyjętego
nominału pieniężnego.

Brak modelu salda oznacza, że środowisko nie bada między innymi:

\begin{itemize}
    \item ryzyka bankructwa,
    \item doboru wysokości podstawowej stawki,
    \item zarządzania kapitałem pomiędzy rozdaniami,
    \item warunków zakończenia całej sesji,
    \item wpływu limitów stołu.
\end{itemize}

Zagadnienia te dotyczą zarządzania kapitałem, a nie samej strategii decyzji
Play. Mogłyby zostać dodane jako zewnętrzna warstwa zarządzająca serią
epizodów.

\subsection{Model pojedynczego gracza}
\label{subsec:model-pojedynczego-gracza}

Środowisko reprezentuje jednego gracza oraz jednego krupiera. Nie symuluje
innych miejsc przy stole ani kart rozdanych pozostałym uczestnikom.

W rzeczywistej grze obecność innych graczy wpływa na liczbę kart zużytych
z talii, ale zakryte karty tych osób pozostają nieznane badanemu graczowi.
Modelowanie pełnego stołu zwiększałoby złożoność procesu rozdawania bez
zmiany informacji dostępnych agentowi w podstawowym problemie decyzyjnym.

Przyjęte uproszczenie oznacza, że wszystkie niewidoczne karty poza kartami
krupiera, kartami spalonymi i odkrytym stołem pozostają w źródle kart.
Rozkład kart widocznych dla pojedynczego gracza jest dzięki temu zgodny
z losowaniem bez zwracania ze standardowej talii.

\subsection{Ograniczenie do jednej konfiguracji reguł}
\label{subsec:jedna-konfiguracja-regul}

Obecna implementacja wykorzystuje jedną konfigurację zasad:

\begin{itemize}
    \item krupier kwalifikuje się od jednej pary,
    \item Ante i Blind mają wartość jednej jednostki,
    \item dostępne mnożniki Play wynoszą \(1\), \(2\), \(3\) i \(4\),
    \item stosowana jest jedna tabela wypłat Blind.
\end{itemize}

Niektóre kasyna lub materiały źródłowe mogą stosować inne tabele wypłat
zakładów dodatkowych albo lokalne warianty zasad. W aktualnym silniku
wartości podstawowej konfiguracji są jawnie zapisane w module reguł
i rozliczeń.

Możliwym rozszerzeniem jest wprowadzenie niemodyfikowalnego obiektu
konfiguracyjnego przekazywanego do silnika. Pozwoliłby on wybierać tabelę
Blind, warunek kwalifikacji oraz aktywne zakłady bez modyfikowania kodu
mechanizmu rozliczeń. Rozszerzenie to nie jest konieczne dla eksperymentów
prowadzonych w jednej, wcześniej ustalonej konfiguracji.

\subsection{Brak numerycznej reprezentacji obserwacji}
\label{subsec:brak-numerycznej-obserwacji}

Obecna obserwacja jest obiektem domenowym zawierającym karty, fazę oraz zbiór
legalnych akcji. Taka postać jest czytelna, jednoznaczna i wygodna podczas
testowania, ale nie może zostać bezpośrednio przekazana do większości modeli
uczenia maszynowego.

Przed rozpoczęciem uczenia konieczne będzie utworzenie encoderów
przekształcających \texttt{UTHObservation} do wektorów lub struktur
numerycznych o stałym rozmiarze. Sposób kodowania nie zostanie umieszczony
w klasie \texttt{UTHGame}, ponieważ porównanie reprezentacji stanu stanowi
jeden z elementów badawczych pracy.

Oddzielenie encodera od silnika umożliwi wykorzystanie identycznych rozdań
i reguł przy porównywaniu różnych reprezentacji. Różnice w wynikach nie będą
wówczas wynikały z odmiennego działania mechanizmu gry.

\section{Podsumowanie}
\label{sec:podsumowanie-srodowiska-uth}

W rozdziale przedstawiono projekt i implementację środowiska pojedynczego
rozdania Ultimate Texas Hold’em. Silnik został podzielony na niezależne
komponenty odpowiedzialne za reprezentację kart, ocenę układów, rozdawanie,
przechowywanie stanu, walidację akcji, przeprowadzanie showdownu oraz
rozliczanie zakładów.

Przebieg rozdania odwzorowano jako maszynę stanów obejmującą fazy
\texttt{PREFLOP}, \texttt{FLOP}, \texttt{RIVER} i \texttt{TERMINAL}.
Zbiór legalnych akcji zależy od bieżącej fazy, a każda akcja typu
\texttt{BET} kończy część decyzyjną rozdania. Konstrukcja ta gwarantuje,
że zakład Play zostanie wykonany najwyżej jeden raz.

Istotnym elementem projektu jest rozdzielenie pełnego stanu
\texttt{RoundState} od obserwacji \texttt{UTHObservation}. Agent otrzymuje
wyłącznie własne karty, odkryte karty wspólne, aktualną fazę oraz legalne
akcje. Karty krupiera, karty spalone i kolejność pozostałej talii pozostają
wewnętrznymi informacjami środowiska. Zapobiega to wykorzystaniu przez
agenta informacji niedostępnych graczowi w rzeczywistym rozdaniu.

Showdown wykorzystuje ogólny evaluator pokerowy wybierający najlepsze pięć
kart z siedmiu. Rozliczenie Ante, Blind i Play jest wykonywane na podstawie
wyniku porównania, kwalifikacji krupiera, mnożnika Play oraz kategorii ręki
gracza. Wszystkie wartości finansowe są przechowywane dokładnie za pomocą
typu \texttt{Fraction}.

Silnik umożliwia zarówno losowe, jak i deterministyczne rozdawanie kart.
Wstrzykiwane źródło \texttt{FixedDeck} pozwala tworzyć powtarzalne scenariusze
testowe obejmujące pełną kolejność kart. Opcjonalny mechanizm historii
zapisuje obserwacje widoczne przed decyzjami oraz końcowy pełny stan
diagnostyczny.

Zaimplementowane środowisko stanowi warstwę domenową projektu. Nie zawiera
jeszcze numerycznego kodowania obserwacji, funkcji nagrody zwracanej przez
interfejs uczenia ani adaptera Gymnasium. Elementy te zostaną utworzone jako
osobne warstwy, dzięki czemu możliwe będzie porównywanie reprezentacji stanu,
funkcji nagrody i algorytmów przy zachowaniu identycznych zasad gry.

Gotowy silnik zapewnia podstawę do implementacji agentów bazowych,
infrastruktury eksperymentalnej oraz metod uczenia przez wzmacnianie
analizowanych w dalszej części pracy.
